% TOP_PROBLEM: {"problemId": "problem.furstenberg-no-bigeodesics-conjecture-planar-fpp", "canonicalTitle": "Furstenberg No-Bigeodesics (Planar FPP, Standard Minimum Moment)", "releaseRank": 485, "primaryDomain": "probability_ergodic_dynamics", "primaryDomainLabel": "Probability, ergodic theory & dynamics", "displayStatus": "open_with_solved_subcases", "releaseStatus": "open_with_solved_subcases"}
% !TeX program = lualatex
% Definition and sources only; this file is not an LLM solution attempt.
\documentclass[11pt]{article}
\usepackage[margin=25mm]{geometry}
\usepackage{fontspec}
\setmainfont{Latin Modern Roman}
\usepackage{amsmath,amssymb,mathtools}
\usepackage{unicode-math}
\setmathfont{Latin Modern Math}
\usepackage{xurl,hyperref}
\hypersetup{colorlinks=true,urlcolor=blue}
\setcounter{secnumdepth}{0}
\setlength{\parindent}{0pt}
\setlength{\parskip}{5pt}
\setlength{\emergencystretch}{3em}
\begin{document}
\section*{\texorpdfstring{Furstenberg No-Bigeodesics (Planar FPP, Standard Minimum Moment)}{Furstenberg No-Bigeodesics (Planar FPP, Standard Minimum Moment)}}\label{485}
\subsection{Definitions and mathematical statement}
Assign independent edge times with a non-atomic law $G$ on $[0,\infty)$ to the unoriented nearest-neighbor edges of ${\mathbb{Z}}^2$. Assume
\[
 {\mathbb{E}}\bigl[\min(t_1,t_2,t_3,t_4)^2\bigr]<\infty
\]
for four independent samples from $G$. For vertices $x,y$, define $T(x,y)$ as the infimum of the sum of edge times over finite paths from $x$ to $y$. A bigeodesic is a doubly infinite self-avoiding nearest-neighbor path $(v_i)_{i\in{\mathbb{Z}}}$ satisfying
\[
 \sum_{k=i}^{j-1}t_{\{v_k,v_{k+1}\}}=T(v_i,v_j)
 \quad\text{for every }i<j.
\]
The assertion is that the product probability of existence of any bigeodesic is zero. The moment condition concerns the minimum of four times, not an individual edge time. No prescribed directions, density for $G$, curvature of a limit shape, or exponential moment is assumed. This selected version is not equated with a separately stated moment-free version.
\par\smallskip\textit{Formulation and concept references:} \href{https://arxiv.org/pdf/1512.00804}{[S1]}.

\subsection{Short English statement}
Under the specified non-atomic minimum-moment assumption, is every doubly infinite path in planar first-passage percolation almost surely nonoptimal somewhere along its length?
\subsection{Sources}
\begin{itemize}
\item[S1]Furstenberg No-Bigeodesics (Planar FPP, Standard Minimum Moment); exact editorial selection using the primary formulations and hypotheses recorded in the source receipts.
\url{https://arxiv.org/pdf/1512.00804}.
\textit{Formulation source.}
\end{itemize}
\end{document}
